\documentclass{article}
\linespread{1.3}
\usepackage[margin=50pt]{geometry}
\usepackage{amsmath, amsthm, amssymb, amsthm, tikz, fancyhdr, float, graphicx, spverbatim, systeme}
\pagestyle{fancy}
\renewcommand{\headrulewidth}{0pt}
\newcommand{\changefont}{\fontsize{15}{15}\selectfont}

\fancypagestyle{firstpageheader}{
  \fancyhead[R]{
    \changefont
    \parbox[t]{4cm}{ % Adjust width as needed
      Michael Huang\\
      EN.555.644.81\\
      Homework 2
    }
  }
}

\begin{document}

\thispagestyle{firstpageheader}
{\Large 

\section*{3.2}

Under what circumstances does a minimum-variance hedge portfolio lead to no hedging at all? \\ \\ 
As the formula implies that \(h = \rho \frac{\sigma_s}{\sigma_F}\), we look for \(h = 0\). As the \(\sigma\) values are realistically unlikely to be 0 since prices tend to naturally fluctuate. That leaves \(\rho = 0\), which indicates that the correlation between the change in the spot price and the change in the futures price is zero.

}
{\Large 

\section*{3.4}

A company has a \$20 million portfolio with a beta of 1.2. It would like to use futures contracts
on a stock index to hedge its risk. The index futures is currently standing at 1080, and each
contract is for delivery of \$250 times the index. What is the hedge that minimizes risk? What
should the company do if it wants to reduce the beta of the portfolio to 0.6?


}
{\Large 

\section*{3.7}

Explain why a short hedger's position improves when the basis strengthens unexpectedly and
worsens when the basis weakens unexpectedly.


}
{\Large 

\section*{3.14}

A corn farmer argues “I do not use futures contracts for hedging. My real risk is not the price of
corn. It is that my whole crop gets wiped out by the weather.”Discuss this viewpoint. Should the
farmer estimate his or her expected production of corn and hedge to try to lock in a price for
expected production?


}
{\Large 

\section*{3.15}

On July 1, an investor holds 50,000 shares of a certain stock. The market price is \$30 per share.
The investor is interested in hedging against movements in the market over the next month and
decides to use the September Mini S\&P 500 futures contract. The index is currently 1,500 and
one contract is for delivery of \$50 times the index. The beta of the stock is 1.3. What strategy
should the investor follow? Under what circumstances will it be profitable?


}
{\Large 

\section*{3.XX}

A futures contract is used for hedging. Explain why the daily settlement of the contract can give
rise to cash flow problems.


}
{\Large 

\section*{3.18}

Suppose the one-year gold lease rate is 1.5% and the one-year risk-free rate is 5.0%. Both rates
are compounded annually. Use discussion in Business Snapshot 3.1 to calculate the maximum
one-year forward price Goldman Sachs should quote for gold when the spot price is \$1,200.


}
{\Large 

\section*{3.XX}

It is now October 2017. A company anticipates that it will purchase 1 million pounds of copper
in each of February 2018, August 2018, February 2019, and August 2019. The company has
decided to use the futures contracts traded by the CME Group to hedge its risk. One contract is
for the delivery of 25,000 pounds of copper. The initial margin is \$2,000 per contract and the
maintenance margin is \$1,500 per contract. The company's policy is to hedge 80\% of its
exposure. Contracts with maturities up to 13 months into the future are considered to have
sufficient liquidity to meet the company's needs. Devise a hedging strategy for the company. (Do
not make the “tailing” adjustment described in Section 3.4.)
Assume the market prices (in cents per pound) today and at future dates are as follows. \\ \\
What is the impact of the strategy you propose on the price the company pays for copper? What
is the initial margin requirement in October 2017? Is the company subject to any margin calls? \\ \\


}

\end{document}